% =====================================================================
% forge :: prototype.tex  (v0.4)
% Complete, compilable example document. The runtime entry point that
% forge clones. Compile with:  latexmk -pdf prototype.tex
% (expects preamble.tex / proofs.tex / skeleton.tex in same directory).
% Sample content uses the endo-blotto example so the scaffold looks real.
% =====================================================================
\documentclass[10pt,a4paper]{article}

% =====================================================================
% forge :: preamble.tex  (v0.4)
% Bundled preamble cloned at the start of every forge run.
% Carries the theorem block + myproof env, the disc design
% (ONE public macro \disc{v|p|h|c} + [v|p|h|c] coloured discs on the
% result environments), per-key AUTO-COUNT totals persisted to .aux,
% the RARITY star strip (\rarity{<0..5 in 0.5 steps>}), and the full-page
% collectible-card cover (\makeprotocover -- v0.4: a standalone first page
% with a STORY block, a tier NAME beside the stars, and a described
% CONFIDENCE legend + RARITY-tier legend).
% pdflatex-safe. Stars via fontawesome5. xcolor arrives via tcolorbox[most].
% =====================================================================

% --- fonts (adopted from main_r1.tex) ---
\usepackage[T1]{fontenc}
\usepackage{textcomp}
\usepackage{libertine}
\usepackage[libertine]{newtxmath}

% --- core math / symbols ---
\usepackage{amsmath}
\usepackage{amsfonts}
\let\Bbbk\relax  % silence duplicate-definition warning (newtxmath vs amssymb)
\usepackage{amssymb}
\usepackage[most]{tcolorbox}  % grey/coloured boxes; also pulls in xcolor
\usepackage{array}            % \extrarowheight + p-columns for the skeleton table
\usepackage{fontawesome5}     % \faStar / \faStarHalf / \faStar[regular] for \rarity

% --- geometry: more whitespace than the old 1in (centred, ~1.3in) ---
\usepackage[a4paper,margin=1.3in]{geometry}

% --- math operator shortcuts (subset from main_r1.tex) ---
\newcommand*{\prob}{\mathsf{P}}  % probability
\newcommand*{\ex}{\mathsf{E}}    % expectation

% =====================================================================
% Theorems and proofs
% =====================================================================
% amsthm is loaded for \newtheoremstyle (the disc design). It does not
% conflict with the custom myproof env below (we do NOT use amsthm's proof).
% newtxmath already defines \openbox, so neutralise it before amsthm
% redefines it (otherwise: "Command \openbox already defined").
\let\openbox\relax
\usepackage{amsthm}

% --- status colours (xcolor available via tcolorbox[most]) ---
\definecolor{stverified}{RGB}{34,139,34}    % green  -> verified
\definecolor{stplausible}{RGB}{31,119,180}  % blue   -> plausible
\definecolor{stheuristic}{RGB}{230,140,0}   % orange -> heuristic
\definecolor{stconjectured}{RGB}{200,30,30} % red    -> conjectured

% --- the disc glyph: a colour-filled bullet, nudged onto the baseline.
% Drawn as a scaled $\bullet$ via \textcolor so NO TikZ dependency is needed.
% Internal helper only; the single PUBLIC macro is \disc{v|p|h|c} below.
\newcommand{\discbullet}[1]{{\textcolor{#1}{\raisebox{0.05ex}{\large$\bullet$}}}}

% --- THE one public macro: \disc{X}, X in {v,p,h,c}. One alphabet everywhere.
%     v = verified (green), p = plausible (blue),
%     h = heuristic (orange), c = conjectured (red).
%     Unknown/empty keys print nothing (graceful).
\newcommand{\disc}[1]{%
  \ifx\\#1\\%                      % empty arg -> nothing
  \else
    \begingroup
    \def\tmp{#1}%
    \def\kv{v}\ifx\tmp\kv\discbullet{stverified}\fi
    \def\kp{p}\ifx\tmp\kp\discbullet{stplausible}\fi
    \def\kh{h}\ifx\tmp\kh\discbullet{stheuristic}\fi
    \def\kc{c}\ifx\tmp\kc\discbullet{stconjectured}\fi
    \endgroup
  \fi
}

% =====================================================================
% Rarity strip :: \rarity{<n>}  (the AMBITION axis -- distinct from POWER)
% =====================================================================
% n in {0,0.5,1,...,5}. Renders five stars out of five via fontawesome5:
%   filled  \faStar          (full point)
%   half    \faStarHalf      (the .5 point)
%   empty   \faStar[regular] (outline -- unearned point)
% Gold-toned to read as a separate axis from the coloured POWER discs.
% pdflatex-safe: tested on TeX Live 2024 (\faStarHalfAlt does NOT exist
% here; \faStarHalf is the correct half-star glyph).
\definecolor{strarity}{RGB}{200,150,0}  % gold -> ambition / rarity
%
% Implementation: work in HALF-STAR UNITS to keep everything integer and
% avoid any float arithmetic. \rar@two = round(2*#1), so each whole star is
% worth 2 units and a half star is 1 unit. Parsing 2*#1 from an argument
% that may carry ".5": \rar@parse strips an optional ".5" tail.
%   "4.5" -> 2*4 + 1 = 9 ;  "3" -> 2*3 = 6.
% NOTE: preamble is \input in the BODY where @ is catcode-12, so every
% @-bearing control sequence here MUST sit inside \makeatletter ... .
\makeatletter
\newcount\rar@two
\newcount\rar@slot
\newcount\rar@rem
% \rar@parse #1 -> sets \rar@two = 2*#1 (handles integer or x.5).
% Trick: append ".0" so the fractional field is ALWAYS a digit, whether the
% caller wrote "3" (-> 3.0.0) or "4.5" (-> 4.5.0). \rar@split keeps the
% integer part #1 and the FIRST fractional digit #2; trailing tokens #3 are
% swallowed up to the \relax sentinel.
\def\rar@parse#1{\edef\rar@arg{#1}\expandafter\rar@split\rar@arg.0.\relax}
\def\rar@split#1.#2#3\relax{%
  \rar@two=#1\relax \multiply\rar@two by 2\relax
  % #2 is the first fractional digit; a leading 5 means a half-star.
  \ifnum#2=5 \advance\rar@two by 1\relax\fi}
%
% \rarity{<n>}: parse to half-units, then walk 5 slots. Slot k (1..5) is
% full if >=2 units remain, half if exactly 1 remains, else empty.
\newcommand{\rarity}[1]{%
  \begingroup
  \rar@parse{#1}%                          \rar@two = 2*n  (half-units)
  {\color{strarity}%
   \rar@slot=0
   \loop
     \advance\rar@slot by 1
     \rar@rem=\rar@two
     \advance\rar@rem by -\numexpr 2*(\rar@slot-1)\relax  % units left for slot
     \ifnum\rar@rem>1 \faStar
     \else\ifnum\rar@rem=1 \faStarHalf
     \else {\faStar[regular]}%
     \fi\fi
   \ifnum\rar@slot<5 \repeat}%
  \endgroup
}
\makeatother
% Four total-counters tick up each time a result env is opened with
% [v|p|h|c]. totcount persists each grand total to the .aux file, so the
% COVER (page 1, before the blocks) prints CORRECT counts after latexmk's
% normal multi-pass run. Exposed as \protoTotV \protoTotP \protoTotH
% \protoTotC, each expanding to the persisted total.
\usepackage{totcount}
\newtotcounter{proto@v}\newtotcounter{proto@p}%
\newtotcounter{proto@h}\newtotcounter{proto@c}

% \protototal{<totcounter>}: robustly expand to the persisted grand total,
% yielding 0 (not an empty token) before the .aux has been read on pass 1.
% This keeps \ifnum\protoTotV>0 well formed on every pass.
\newcommand{\protototal}[1]{%
  \ifcsname c@#1@totc\endcsname \totvalue{#1}\else 0\fi}
\newcommand{\protoTotV}{\protototal{proto@v}}
\newcommand{\protoTotP}{\protototal{proto@p}}
\newcommand{\protoTotH}{\protototal{proto@h}}
\newcommand{\protoTotC}{\protototal{proto@c}}

% \protocount{<key>}: bump the matching total-counter. Called by \dischead.
\newcommand{\protocount}[1]{%
  \def\tmp{#1}%
  \def\kv{v}\ifx\tmp\kv\stepcounter{proto@v}\fi
  \def\kp{p}\ifx\tmp\kp\stepcounter{proto@p}\fi
  \def\kh{h}\ifx\tmp\kh\stepcounter{proto@h}\fi
  \def\kc{c}\ifx\tmp\kc\stepcounter{proto@c}\fi
}

% --- head wrapper: bump the count, then render disc + thin space, but ONLY
%     when the key is non-empty (an empty slot leaves no stray space).
\newcommand{\dischead}[1]{%
  \ifx\\#1\\\else\protocount{#1}\disc{#1}\thinspace\fi
}

% --- custom theorem style: NOTE (#3 = the [v|p|h|c] key) rendered FIRST as a
%     disc, then name+number. amsthm delivers \begin{env}[<key>] into #3 free.
\newtheoremstyle{status}%
  {\topsep}%                                  space above
  {\topsep}%                                  space below
  {\itshape}%                                 body font
  {}%                                         indent
  {\bfseries}%                                head font
  {.}%                                        punctuation after head
  {.5em}%                                     space after head
  {\dischead{#3}\thmname{#1}\thmnumber{ #2}}% HEAD: disc{#3} BEFORE name+number
\theoremstyle{status}

% --- the eight environments (same names/labels as main_r1.tex) ---
\newtheorem{claim}{Claim}
\newtheorem{assumption}{Assumption}
\newtheorem{theorem}{Theorem}
\newtheorem{proposition}{Proposition}
\newtheorem{definition}{Definition}
\newtheorem{lemma}{Lemma}
\newtheorem{corollary}{Corollary}
\newtheorem{remark}{Remark}

% --- custom proof environment (verbatim from main_r1.tex, line 76).
%     #1 = type label (e.g. Proposition), #2 = the \ref{...}.
%     Kept hand-rolled; we do NOT use amsthm's proof.
\newenvironment{myproof}[2]{\paragraph{Proof of {#1} {#2}.}}{\hfill $\blacksquare$}

% =====================================================================
% Building-block visual furniture
% =====================================================================
% \blocktag{<text>}: a thin full-width rule + a small bold tag, so each
% building block reads as a distinct card rather than a wall of text.
\newcommand{\blocktag}[1]{%
  \par\bigskip
  {\color{gray!55}\rule{\linewidth}{0.4pt}}\par
  \nobreak\vspace{2pt}%
  {\footnotesize\bfseries\color{gray!70}#1}\par
  \nobreak\vspace{2pt}%
}

% Inline part labels for inside a block.
\newcommand{\Intuition}{\noindent\textbf{Intuition.}\ }
\newcommand{\Proofsketch}{\noindent\textbf{Proof sketch.}\ }
\newcommand{\Risk}{\noindent\textbf{Risk.}\ }
\newcommand{\Uses}{\noindent\textbf{Uses $\rightarrow$}\ }

% =====================================================================
% Collectible-card cover :: \makeprotocover   (v0.4 -- full first page)
% =====================================================================
% Reads the document-set macros (all defined before \begin{document}):
%   \protonum \protoslug \protofield \protodate \prototitle
%   \protosubtitle \protostory \protopunch \protospine
%   \protorarity (0..5 in .5 steps) \prototier (a tier NAME, never a journal)
%   \protoraritynote (one short line, tier language only)
% The card is its OWN full first page (titlepage-style, \thispagestyle{empty}),
% then \clearpage hands a fresh page 2 to the abstract/body.
% Two distinct axes:
%   * RARITY = gold stars (\rarity), AMBITION, with a tier NAME alongside.
%   * CONFIDENCE = coloured POWER discs (\protoTotV/P/H/C, auto-counted).
% The totals expand to \value{...} which is NOT digestible by \ifnum, so we
% stage each into a scratch counter and read it back number-safely.
\newcounter{proto@scratch}
\newcommand{\makeprotocover}{%
\thispagestyle{empty}%
\begingroup
\setlength{\parindent}{0pt}%
\null\vspace*{\fill}
\begin{tcolorbox}[enhanced, breakable=false, sharp corners=downhill,
                  colback=gray!2, colframe=gray!55, boxrule=0.8pt,
                  arc=4pt, left=16pt, right=16pt, top=14pt, bottom=14pt,
                  drop shadow southeast]
  % --- top row: id / field / date  ---  forge stamp ---
  {\footnotesize\color{gray!75}%
    \textbf{№\,\protonum}\;\textperiodcentered\;\protofield
    \;\textperiodcentered\;\protodate
    \hfill \textsc{forge \textperiodcentered\ beta 0.1}}\par
  \vspace{5pt}
  {\color{gray!40}\rule{\linewidth}{0.4pt}}\par
  \vspace{10pt}
  % --- title / subtitle ---
  {\LARGE\bfseries \prototitle\par}
  \vspace{3pt}
  {\itshape\color{gray!75} \protosubtitle\par}
  \vspace{12pt}
  % --- STORY block: the most important line on the card ---
  {\footnotesize\bfseries\color{gray!70}STORY\par}
  \vspace{3pt}
  {\protostory\par}
  \vspace{8pt}
  % --- punch line (the one-line mechanism / contribution) ---
  {\color{stplausible}$\blacktriangleright$}\ \protopunch\par
  \vspace{12pt}
  {\color{gray!40}\rule{\linewidth}{0.4pt}}\par
  \vspace{9pt}
  % --- rarity line: stars (n/5) + tier NAME, then the note in gray italic ---
  {\textbf{RARITY.}\ %
    \rarity{\protorarity}\ %
    {\color{strarity}(\protorarity/5)}\;\textperiodcentered\;%
    \textsc{\prototier}\par}
  \vspace{2pt}
  {\footnotesize\itshape\color{gray!70}\protoraritynote\par}
  \vspace{9pt}
  % --- spine ---
  {\textbf{SPINE.}\ \protospine\par}
  \vspace{9pt}
  {\color{gray!40}\rule{\linewidth}{0.4pt}}\par
  \vspace{9pt}
  % --- CONFIDENCE legend: described, one level per line, with live tally ---
  {\footnotesize\bfseries\color{gray!70}CONFIDENCE \textnormal{\color{gray!60}%
     (coloured discs --- how sure each result is)}\par}
  \vspace{3pt}
  {\footnotesize
   \setcounter{proto@scratch}{\protoTotV}%
     \disc{v}~\textbf{Verified}\ \textbf{(\arabic{proto@scratch})} ---
     a complete proof checked, or numerically confirmed by \textsc{temper}\par
   \setcounter{proto@scratch}{\protoTotP}%
     \disc{p}~\textbf{Plausible}\ \textbf{(\arabic{proto@scratch})} ---
     the key step is checked; the proof is not yet complete\par
   \setcounter{proto@scratch}{\protoTotH}%
     \disc{h}~\textbf{Heuristic}\ \textbf{(\arabic{proto@scratch})} ---
     intuition or analogy only; unchecked\par
   \setcounter{proto@scratch}{\protoTotC}%
     \disc{c}~\textbf{Conjectured}\ \textbf{(\arabic{proto@scratch})} ---
     believed; no argument yet\par}
  \vspace{4pt}
  {\footnotesize\color{gray!75}%
    \setcounter{proto@scratch}{\protoTotV}%
    \ifnum\value{proto@scratch}>0
      \disc{v}\ \arabic{proto@scratch}\ verified so far.%
    \else
      \disc{v}\ 0 verified --- run \textsc{temper} to charge them green.%
    \fi\par}
  \vspace{9pt}
  {\color{gray!40}\rule{\linewidth}{0.4pt}}\par
  \vspace{9pt}
  % --- RARITY-tier legend (gold stars = ambition; distinct from confidence) ---
  {\footnotesize\bfseries\color{gray!70}RARITY TIERS \textnormal{\color{gray!60}%
     (gold stars --- how ambitious the target is)}\par}
  \vspace{3pt}
  {\footnotesize\color{gray!75}%
    {\color{strarity}5}~top-5 \quad
    {\color{strarity}4}~flagship field \quad
    {\color{strarity}3}~strong field \quad
    {\color{strarity}2}~seed \quad
    {\color{strarity}1}~below-bar\par}
\end{tcolorbox}%
\vspace*{\fill}
\endgroup
\clearpage
}

% =====================================================================
% Footer stamp
% =====================================================================
\newcommand{\forgestamp}{prototype \textperiodcentered\ forge \textperiodcentered\ beta 0.1}
\usepackage{fancyhdr}
\pagestyle{fancy}
\fancyhf{}
\renewcommand{\headrulewidth}{0pt}
\renewcommand{\footrulewidth}{0pt}
\fancyfoot[C]{\footnotesize\color{gray!70}\forgestamp}
\fancyfoot[R]{\footnotesize\color{gray!70}\thepage}


% --- document-set macros read by \makeprotocover (define BEFORE document) ---
\newcommand{\protonum}{007}
\newcommand{\protoslug}{endo-blotto}
\newcommand{\protofield}{econ}
\newcommand{\protodate}{\today}
\newcommand{\prototitle}{Endogenous Battlefields in Colonel Blotto}
\newcommand{\protosubtitle}{When contestants choose where to fight before how hard}
% --- STORY: the overall world of the paper -- who needs the result and why.
%     DIFFERENT from \protopunch (the one-line mechanism/contribution below). ---
\newcommand{\protostory}{Defenders often must build up what they are about to
  fight over --- infrastructure, capacity, reputation --- before the contest
  begins. This paper asks how that pre-conflict investment should be spread when
  the thing being protected has diminishing value, and shows the answer flips the
  classic ``concentrate everything'' prescription.}
\newcommand{\protopunch}{Letting players pick the battlefield set first collapses the
  classic mixed equilibrium onto a sparse, characterisable support.}
\newcommand{\protospine}{Lemma~\ref{lem:concentrate} $\to$ Corollary~\ref{cor:sparse}
  $\to$ Proposition~\ref{prop:exist} $\Rightarrow$ Theorem~\ref{thm:char}}
% --- rarity (AMBITION axis): n in {0,0.5,...,5}, a tier NAME, + a one-line why.
%     \prototier is a GENERAL tier name only -- NEVER a specific journal. ---
\newcommand{\protorarity}{3}
\newcommand{\prototier}{strong field}
\newcommand{\protoraritynote}{a real seed today; strong-field ceiling if the
  global-curvature claim closes.}
% --- RARITY RATIONALE: 1-2 sentences, rendered at the END, on WHY this tier ---
%     what earns the rung, and what would lift or lower it. Fuller than the
%     card's one-line \protoraritynote. ---
\newcommand{\protorarityrationale}{Strong-field, not flagship: it puts a genuinely
  new, tractable twist---endogenous front-selection---on a canonical model and
  extracts a clean closed form, which is novel enough to place well. But it is one
  mechanism on a known game rather than a new framework; it would climb a tier if
  the characterisation extended past the symmetric case or anchored a concrete
  application.}

\title{\prototitle}
\date{\protodate}

\begin{document}

% The cover owns page 1 (\thispagestyle{empty}) and \clearpage's to page 2,
% where \pagestyle{fancy} (set in the preamble) resumes for the body.
\makeprotocover

\begin{abstract}
\noindent We study Colonel Blotto with an \emph{endogenous} battlefield stage:
each player first selects a subset of fronts to open, then allocates a fixed
budget across the union of opened fronts. The selection stage is shown to
concentrate effort, which sparsifies the support of any equilibrium allocation
and lets us characterise the resulting mixed equilibrium in closed form. This is
a forge prototype: every result carries a coloured status disc, an intuition
line, a footnote-sized proof sketch, and a forward dependency pointer. Proof
sketches sit in the appendix; the dependency-flow skeleton closes the document.
\end{abstract}

% =====================================================================
\section{Setup}\label{sec:setup}
% =====================================================================

\paragraph{Players} \noindent Two players $i\in\{A,B\}$, each holding a budget
normalised to $1$ and facing $n$ potential fronts.

\paragraph{Timing} \noindent The game proceeds in two stages:
\begin{enumerate}\setlength{\itemsep}{1pt}
  \item Selection stage: each player picks a set $S_i\subseteq\{1,\dots,n\}$ of
        fronts to open, at cost $\kappa$ per front.
  \item Allocation stage: players simultaneously allocate their unit budget
        across the contested union $S_A\cup S_B$, winning a front by outspending
        the rival there.
\end{enumerate}

\paragraph{Strategies} \noindent A pure strategy is a pair: an open set
$S_i\subseteq\{1,\dots,n\}$ and a budget vector on $S_A\cup S_B$ summing to $1$.
Mixed strategies are distributions over such pairs; the second-stage allocation
is mixed in equilibrium.

\paragraph{Payoffs} \noindent A player's payoff is the number of fronts won
minus opening costs: front $f$ is won by whoever spends strictly more there
(ties split evenly), and each opened front carries cost $\kappa$.

\paragraph{Solution concept} \noindent Subgame-perfect equilibrium, solved by
backward induction: the allocation stage is fixed first for each $(S_A,S_B)$,
then the selection stage is optimised against it.

\begin{assumption}[v]\label{ass:cost}
The per-front opening cost satisfies $0<\kappa<1/n$, and ties on a front are
split evenly.
\end{assumption}

\noindent\Intuition The cost is small enough that opening at least one front is
always worthwhile, but positive, so players never open fronts they cannot
contest.

\begin{proposition}[p]\label{prop:exist}
For any fixed pair of open sets $(S_A,S_B)$, the second-stage allocation game
has a mixed equilibrium.
\end{proposition}

\noindent\Intuition The reduced allocation game is a finite-budget Blotto on the
contested union; existence follows from the standard concavification of the
per-front contest, independently of how the fronts were selected.

% =====================================================================
\section{Building blocks}\label{sec:blocks}
% =====================================================================

\blocktag{BLOCK 1 \textperiodcentered\ concentration lemma}
\begin{lemma}[p]\label{lem:concentrate}
In any equilibrium, each player opens at most $\lceil 1/(2\kappa)\rceil$ fronts.
\end{lemma}

\Intuition Spreading the budget thinly makes every front cheap to steal, so a
deviation that abandons a thin front and reinforces another strictly pays once
the opening cost is sunk; the cap is the point where this stops paying.

{\footnotesize\Proofsketch Bound the marginal win-probability of a front
receiving budget $b$ and compare against $\kappa$; fronts below the threshold are
dropped. Check \checkmark\ (symbolic): at $\kappa\to 1/n$ the cap $\to\lceil
n/2\rceil$, matching the all-fronts-open benchmark.\par}

{\footnotesize\Risk The bound assumes a single binding deviation; coalitional or
multi-front deviations are not yet ruled out.\par}

\Uses Assumption~\ref{ass:cost}.

\blocktag{BLOCK 2 \textperiodcentered\ sparse support}
\begin{corollary}[h]\label{cor:sparse}
The equilibrium allocation has support on $O(1/\kappa)$ fronts, independent of
$n$.
\end{corollary}

\Intuition If both players open boundedly many fronts, the contested union is
bounded too, so almost all of the $n$ fronts are never touched and drop out of
the support.

{\footnotesize\Proofsketch Union the two caps from
Lemma~\ref{lem:concentrate}. Check \checkmark\ (symbolic): support size
$\le 2\lceil 1/(2\kappa)\rceil = O(1/\kappa)$, free of $n$.\par}

{\footnotesize\Risk Treats the two players' open sets as independent; overlap
could shrink \emph{or} reshape the support in ways not captured here.\par}

\Uses Lemma~\ref{lem:concentrate}.

\blocktag{BLOCK 3 \textperiodcentered\ characterisation}
\begin{theorem}[c]\label{thm:char}
A symmetric mixed equilibrium exists in which each player opens exactly
$m^\star=\lceil 1/(2\kappa)\rceil$ fronts and allocates uniformly on
$[0,2/m^\star]$ per opened front.
\end{theorem}

\Intuition With a sparse, symmetric support the second stage reduces to a
standard Blotto on $m^\star$ symmetric fronts, whose equilibrium is the uniform
marginal --- the selection stage just fixes which fronts that runs on.

{\footnotesize\Proofsketch Existence of the reduced game by
Proposition~\ref{prop:exist}; plug the sparse support from
Corollary~\ref{cor:sparse} into the uniform-marginal solution. Check \checkmark\
(symbolic): per-front mean $1/m^\star$ sums to budget $1$.\par}

{\footnotesize\Risk Symmetry is imposed, not derived; asymmetric equilibria with
the same support may coexist and are conjectural.\par}

\Uses Corollary~\ref{cor:sparse}, Proposition~\ref{prop:exist}.

% =====================================================================
\section*{Open joints}
% =====================================================================
\begin{itemize}\setlength{\itemsep}{1pt}
  \item Rule out multi-front and coalitional deviations in
        Lemma~\ref{lem:concentrate} (would promote \disc{p}$\to$\disc{v}).
  \item Establish uniqueness, not just existence, of the symmetric equilibrium
        in Theorem~\ref{thm:char} (currently \disc{c}).
  \item Handle the overlap of $S_A\cap S_B$ explicitly in
        Corollary~\ref{cor:sparse}.
\end{itemize}

% =====================================================================
\appendix
% =====================================================================
% forge :: proofs.tex  (v0.2)
% Proof tier (Appendix A). myproof sketches for the building blocks in
% prototype.tex. Each myproof takes #1 = env type, #2 = \ref{...}.
% Whole section is \footnotesize. % =====================================================================
% forge :: proofs.tex  (v0.2)
% Proof tier (Appendix A). myproof sketches for the building blocks in
% prototype.tex. Each myproof takes #1 = env type, #2 = \ref{...}.
% Whole section is \footnotesize. % =====================================================================
% forge :: proofs.tex  (v0.2)
% Proof tier (Appendix A). myproof sketches for the building blocks in
% prototype.tex. Each myproof takes #1 = env type, #2 = \ref{...}.
% Whole section is \footnotesize. \input{proofs} after \appendix.
% =====================================================================
\section{Proof sketches}\label{app:proofs}

\footnotesize

\begin{myproof}{Proposition}{\ref{prop:exist}}
Fix $(S_A,S_B)$ and let $F=S_A\cup S_B$ be the contested fronts. Each player's
strategy is a distribution over budget vectors on $F$ summing to $1$. The
per-front win payoff is upper semicontinuous in own spend and the strategy sets
are compact and convex after the usual mixing, so a fixed point of the
best-response correspondence exists by Glicksberg. Check \checkmark\ (symbolic):
for $|F|=1$ the game collapses to ``spend everything,'' a degenerate equilibrium,
consistent with the general statement.
\end{myproof}

\begin{myproof}{Lemma}{\ref{lem:concentrate}}
Suppose a player opens $k$ fronts and splits the budget so some front receives
$b<\kappa$. The marginal win probability contributed by that front is at most
$b$ (it is outspent with probability $\ge 1-b$), while keeping it open costs
$\kappa>b$. Dropping the front and reinforcing another raises payoff, so no
equilibrium has such a front. Iterating, every open front receives $\ge\kappa$,
forcing $k\kappa\le$ (effective spend) and hence
$k\le\lceil 1/(2\kappa)\rceil$. Check \checkmark\ (symbolic): the
threshold $b=\kappa$ equates marginal gain and cost, the standard first-order
balance.
\end{myproof}

\begin{myproof}{Corollary}{\ref{cor:sparse}}
Each player opens at most $m^\star=\lceil 1/(2\kappa)\rceil$ fronts by
Lemma~\ref{lem:concentrate}, so the contested union $S_A\cup S_B$ has at most
$2m^\star$ fronts; all other fronts receive zero spend from both players and are
not in the support. Thus the support size is $\le 2\lceil 1/(2\kappa)\rceil =
O(1/\kappa)$, independent of $n$. Check \checkmark\ (symbolic): substituting
$\kappa=1/(2m)$ gives support $\le 2m$, exactly the union of two $m$-front sets.
\end{myproof}

\begin{myproof}{Theorem}{\ref{thm:char}}
Restrict to symmetric profiles. By Corollary~\ref{cor:sparse} the support is
sparse, and by Proposition~\ref{prop:exist} the reduced game on any fixed
symmetric $m^\star$-front union has an equilibrium. On symmetric fronts the
unique symmetric second-stage solution is the uniform marginal on
$[0,2/m^\star]$ per front (the classic Blotto marginal), whose mean
$1/m^\star$ summed over $m^\star$ fronts exhausts the budget. The selection
stage is then indifferent across which $m^\star$ fronts to open, pinning
$m^\star=\lceil 1/(2\kappa)\rceil$. Check \checkmark\ (symbolic): per-front mean
$\int_0^{2/m^\star}\!x\,(m^\star/2)\,dx = 1/m^\star$, so total spend $=1$.
(Conjectural: uniqueness and asymmetric equilibria are not established here.)
\end{myproof}
 after \appendix.
% =====================================================================
\section{Proof sketches}\label{app:proofs}

\footnotesize

\begin{myproof}{Proposition}{\ref{prop:exist}}
Fix $(S_A,S_B)$ and let $F=S_A\cup S_B$ be the contested fronts. Each player's
strategy is a distribution over budget vectors on $F$ summing to $1$. The
per-front win payoff is upper semicontinuous in own spend and the strategy sets
are compact and convex after the usual mixing, so a fixed point of the
best-response correspondence exists by Glicksberg. Check \checkmark\ (symbolic):
for $|F|=1$ the game collapses to ``spend everything,'' a degenerate equilibrium,
consistent with the general statement.
\end{myproof}

\begin{myproof}{Lemma}{\ref{lem:concentrate}}
Suppose a player opens $k$ fronts and splits the budget so some front receives
$b<\kappa$. The marginal win probability contributed by that front is at most
$b$ (it is outspent with probability $\ge 1-b$), while keeping it open costs
$\kappa>b$. Dropping the front and reinforcing another raises payoff, so no
equilibrium has such a front. Iterating, every open front receives $\ge\kappa$,
forcing $k\kappa\le$ (effective spend) and hence
$k\le\lceil 1/(2\kappa)\rceil$. Check \checkmark\ (symbolic): the
threshold $b=\kappa$ equates marginal gain and cost, the standard first-order
balance.
\end{myproof}

\begin{myproof}{Corollary}{\ref{cor:sparse}}
Each player opens at most $m^\star=\lceil 1/(2\kappa)\rceil$ fronts by
Lemma~\ref{lem:concentrate}, so the contested union $S_A\cup S_B$ has at most
$2m^\star$ fronts; all other fronts receive zero spend from both players and are
not in the support. Thus the support size is $\le 2\lceil 1/(2\kappa)\rceil =
O(1/\kappa)$, independent of $n$. Check \checkmark\ (symbolic): substituting
$\kappa=1/(2m)$ gives support $\le 2m$, exactly the union of two $m$-front sets.
\end{myproof}

\begin{myproof}{Theorem}{\ref{thm:char}}
Restrict to symmetric profiles. By Corollary~\ref{cor:sparse} the support is
sparse, and by Proposition~\ref{prop:exist} the reduced game on any fixed
symmetric $m^\star$-front union has an equilibrium. On symmetric fronts the
unique symmetric second-stage solution is the uniform marginal on
$[0,2/m^\star]$ per front (the classic Blotto marginal), whose mean
$1/m^\star$ summed over $m^\star$ fronts exhausts the budget. The selection
stage is then indifferent across which $m^\star$ fronts to open, pinning
$m^\star=\lceil 1/(2\kappa)\rceil$. Check \checkmark\ (symbolic): per-front mean
$\int_0^{2/m^\star}\!x\,(m^\star/2)\,dx = 1/m^\star$, so total spend $=1$.
(Conjectural: uniqueness and asymmetric equilibria are not established here.)
\end{myproof}
 after \appendix.
% =====================================================================
\section{Proof sketches}\label{app:proofs}

\footnotesize

\begin{myproof}{Proposition}{\ref{prop:exist}}
Fix $(S_A,S_B)$ and let $F=S_A\cup S_B$ be the contested fronts. Each player's
strategy is a distribution over budget vectors on $F$ summing to $1$. The
per-front win payoff is upper semicontinuous in own spend and the strategy sets
are compact and convex after the usual mixing, so a fixed point of the
best-response correspondence exists by Glicksberg. Check \checkmark\ (symbolic):
for $|F|=1$ the game collapses to ``spend everything,'' a degenerate equilibrium,
consistent with the general statement.
\end{myproof}

\begin{myproof}{Lemma}{\ref{lem:concentrate}}
Suppose a player opens $k$ fronts and splits the budget so some front receives
$b<\kappa$. The marginal win probability contributed by that front is at most
$b$ (it is outspent with probability $\ge 1-b$), while keeping it open costs
$\kappa>b$. Dropping the front and reinforcing another raises payoff, so no
equilibrium has such a front. Iterating, every open front receives $\ge\kappa$,
forcing $k\kappa\le$ (effective spend) and hence
$k\le\lceil 1/(2\kappa)\rceil$. Check \checkmark\ (symbolic): the
threshold $b=\kappa$ equates marginal gain and cost, the standard first-order
balance.
\end{myproof}

\begin{myproof}{Corollary}{\ref{cor:sparse}}
Each player opens at most $m^\star=\lceil 1/(2\kappa)\rceil$ fronts by
Lemma~\ref{lem:concentrate}, so the contested union $S_A\cup S_B$ has at most
$2m^\star$ fronts; all other fronts receive zero spend from both players and are
not in the support. Thus the support size is $\le 2\lceil 1/(2\kappa)\rceil =
O(1/\kappa)$, independent of $n$. Check \checkmark\ (symbolic): substituting
$\kappa=1/(2m)$ gives support $\le 2m$, exactly the union of two $m$-front sets.
\end{myproof}

\begin{myproof}{Theorem}{\ref{thm:char}}
Restrict to symmetric profiles. By Corollary~\ref{cor:sparse} the support is
sparse, and by Proposition~\ref{prop:exist} the reduced game on any fixed
symmetric $m^\star$-front union has an equilibrium. On symmetric fronts the
unique symmetric second-stage solution is the uniform marginal on
$[0,2/m^\star]$ per front (the classic Blotto marginal), whose mean
$1/m^\star$ summed over $m^\star$ fronts exhausts the budget. The selection
stage is then indifferent across which $m^\star$ fronts to open, pinning
$m^\star=\lceil 1/(2\kappa)\rceil$. Check \checkmark\ (symbolic): per-front mean
$\int_0^{2/m^\star}\!x\,(m^\star/2)\,dx = 1/m^\star$, so total spend $=1$.
(Conjectural: uniqueness and asymmetric equilibria are not established here.)
\end{myproof}


% =====================================================================
% Dependency-flow skeleton (vertical map of the argument)
% =====================================================================
% forge :: skeleton.tex  (v0.3)
% Vertical dependency-flow map as a READABLE LOGIC STORY. One page.
% Read down; each step follows from the ones above. Each NODE is a small
% two-line brick:
%   line 1:  <disc>  <FULL object name, bold via \ref>  --  <plain-English role>
%   line 2:  (indented, smaller) <one concrete detail>  [optional (needs <Name>)]
% Between steps: a centered down-arrow carrying the logic word ("so"/"then"/
% "still open"). Dependencies are implied by vertical position; an explicit
% "(needs <Full Name>)" is added ONLY for a cross-link that is not the node
% directly above. Numbers stay synced via \ref to the labels in prototype.tex.
% % =====================================================================
% forge :: skeleton.tex  (v0.3)
% Vertical dependency-flow map as a READABLE LOGIC STORY. One page.
% Read down; each step follows from the ones above. Each NODE is a small
% two-line brick:
%   line 1:  <disc>  <FULL object name, bold via \ref>  --  <plain-English role>
%   line 2:  (indented, smaller) <one concrete detail>  [optional (needs <Name>)]
% Between steps: a centered down-arrow carrying the logic word ("so"/"then"/
% "still open"). Dependencies are implied by vertical position; an explicit
% "(needs <Full Name>)" is added ONLY for a cross-link that is not the node
% directly above. Numbers stay synced via \ref to the labels in prototype.tex.
% % =====================================================================
% forge :: skeleton.tex  (v0.3)
% Vertical dependency-flow map as a READABLE LOGIC STORY. One page.
% Read down; each step follows from the ones above. Each NODE is a small
% two-line brick:
%   line 1:  <disc>  <FULL object name, bold via \ref>  --  <plain-English role>
%   line 2:  (indented, smaller) <one concrete detail>  [optional (needs <Name>)]
% Between steps: a centered down-arrow carrying the logic word ("so"/"then"/
% "still open"). Dependencies are implied by vertical position; an explicit
% "(needs <Full Name>)" is added ONLY for a cross-link that is not the node
% directly above. Numbers stay synced via \ref to the labels in prototype.tex.
% \input{skeleton} from prototype.tex (after the appendix).
% =====================================================================
\clearpage
\section*{Dependency-flow skeleton}
\addcontentsline{toc}{section}{Dependency-flow skeleton}

% A NODE brick: disc, full name (bold), plain-English role; then an indented
% smaller concrete-detail line. #1 disc | #2 \ref name | #3 role | #4 detail.
\newcommand{\sknode}[4]{%
  \noindent #1\, \textbf{#2} \;---\; #3\par
  \nobreak\smallskip
  {\footnotesize\color{gray!75}\hspace*{1.6em}#4\par}}

% The logic arrow between steps, carrying the connective word.
\newcommand{\skstep}[1]{%
  \par\medskip\centerline{$\downarrow$\ \footnotesize\itshape\color{gray!70}#1}\par\medskip}

% A faint right-aligned tier tag (optional, must not dominate).
\newcommand{\sktag}[1]{\hfill{\scriptsize\color{gray!45}#1}}

\begin{center}
\begin{tcolorbox}[enhanced, breakable, width=0.92\linewidth,
                  colback=gray!3, colframe=gray!45, boxrule=0.5pt,
                  arc=3pt, left=12pt, right=12pt, top=8pt, bottom=8pt,
                  title={Argument flow \quad
                         \disc{v}\,V \;\disc{p}\,P \;\disc{h}\,H \;\disc{c}\,C}]

\sknode{\disc{v}}{Assumption~\ref{ass:cost}}{entry costs a little, and ties split evenly}{%
  $0<\kappa<1/n$ per front opened --- worth opening one, never worth opening all\sktag{setup}}

\skstep{so}

\sknode{\disc{p}}{Proposition~\ref{prop:exist}}{the contest has a mixed-strategy equilibrium}{%
  for any opened fronts, players randomize how hard to fight on each\sktag{engine}}

\skstep{and}

\sknode{\disc{p}}{Lemma~\ref{lem:concentrate}}{each side opens only a few fronts}{%
  at most $\lceil 1/(2\kappa)\rceil$ of them --- a bounded number, not all $n$\sktag{engine}}

\skstep{so}

\sknode{\disc{h}}{Corollary~\ref{cor:sparse}}{those few fronts don't grow with $n$}{%
  the support stays size $O(1/\kappa)$, independent of $n$\sktag{hinge}}

\skstep{so}

\sknode{\disc{c}}{Theorem~\ref{thm:char}}{exact closed-form for how they play}{%
  a symmetric mixed equilibrium: open $m^\star=\lceil 1/(2\kappa)\rceil$ fronts,
  spend uniformly on $[0,2/m^\star]$ each
  \;\textnormal{(needs Proposition~\ref{prop:exist})}\sktag{main result}}

\skstep{still open}

\sknode{\disc{c}}{Uniqueness?}{is this the only solution?}{%
  rule out asymmetric equilibria on the same support --- would turn
  Theorem~\ref{thm:char} green\sktag{open frontier}}

\medskip
\hrule
\smallskip
{\footnotesize\itshape
Read top-to-bottom: each step follows from the ones above; the word on each
$\downarrow$ is the logic. ``(needs $X$)'' flags a step that leans on something
not directly above it. Disc colour = current verification status.}

\end{tcolorbox}
\end{center}
 from prototype.tex (after the appendix).
% =====================================================================
\clearpage
\section*{Dependency-flow skeleton}
\addcontentsline{toc}{section}{Dependency-flow skeleton}

% A NODE brick: disc, full name (bold), plain-English role; then an indented
% smaller concrete-detail line. #1 disc | #2 \ref name | #3 role | #4 detail.
\newcommand{\sknode}[4]{%
  \noindent #1\, \textbf{#2} \;---\; #3\par
  \nobreak\smallskip
  {\footnotesize\color{gray!75}\hspace*{1.6em}#4\par}}

% The logic arrow between steps, carrying the connective word.
\newcommand{\skstep}[1]{%
  \par\medskip\centerline{$\downarrow$\ \footnotesize\itshape\color{gray!70}#1}\par\medskip}

% A faint right-aligned tier tag (optional, must not dominate).
\newcommand{\sktag}[1]{\hfill{\scriptsize\color{gray!45}#1}}

\begin{center}
\begin{tcolorbox}[enhanced, breakable, width=0.92\linewidth,
                  colback=gray!3, colframe=gray!45, boxrule=0.5pt,
                  arc=3pt, left=12pt, right=12pt, top=8pt, bottom=8pt,
                  title={Argument flow \quad
                         \disc{v}\,V \;\disc{p}\,P \;\disc{h}\,H \;\disc{c}\,C}]

\sknode{\disc{v}}{Assumption~\ref{ass:cost}}{entry costs a little, and ties split evenly}{%
  $0<\kappa<1/n$ per front opened --- worth opening one, never worth opening all\sktag{setup}}

\skstep{so}

\sknode{\disc{p}}{Proposition~\ref{prop:exist}}{the contest has a mixed-strategy equilibrium}{%
  for any opened fronts, players randomize how hard to fight on each\sktag{engine}}

\skstep{and}

\sknode{\disc{p}}{Lemma~\ref{lem:concentrate}}{each side opens only a few fronts}{%
  at most $\lceil 1/(2\kappa)\rceil$ of them --- a bounded number, not all $n$\sktag{engine}}

\skstep{so}

\sknode{\disc{h}}{Corollary~\ref{cor:sparse}}{those few fronts don't grow with $n$}{%
  the support stays size $O(1/\kappa)$, independent of $n$\sktag{hinge}}

\skstep{so}

\sknode{\disc{c}}{Theorem~\ref{thm:char}}{exact closed-form for how they play}{%
  a symmetric mixed equilibrium: open $m^\star=\lceil 1/(2\kappa)\rceil$ fronts,
  spend uniformly on $[0,2/m^\star]$ each
  \;\textnormal{(needs Proposition~\ref{prop:exist})}\sktag{main result}}

\skstep{still open}

\sknode{\disc{c}}{Uniqueness?}{is this the only solution?}{%
  rule out asymmetric equilibria on the same support --- would turn
  Theorem~\ref{thm:char} green\sktag{open frontier}}

\medskip
\hrule
\smallskip
{\footnotesize\itshape
Read top-to-bottom: each step follows from the ones above; the word on each
$\downarrow$ is the logic. ``(needs $X$)'' flags a step that leans on something
not directly above it. Disc colour = current verification status.}

\end{tcolorbox}
\end{center}
 from prototype.tex (after the appendix).
% =====================================================================
\clearpage
\section*{Dependency-flow skeleton}
\addcontentsline{toc}{section}{Dependency-flow skeleton}

% A NODE brick: disc, full name (bold), plain-English role; then an indented
% smaller concrete-detail line. #1 disc | #2 \ref name | #3 role | #4 detail.
\newcommand{\sknode}[4]{%
  \noindent #1\, \textbf{#2} \;---\; #3\par
  \nobreak\smallskip
  {\footnotesize\color{gray!75}\hspace*{1.6em}#4\par}}

% The logic arrow between steps, carrying the connective word.
\newcommand{\skstep}[1]{%
  \par\medskip\centerline{$\downarrow$\ \footnotesize\itshape\color{gray!70}#1}\par\medskip}

% A faint right-aligned tier tag (optional, must not dominate).
\newcommand{\sktag}[1]{\hfill{\scriptsize\color{gray!45}#1}}

\begin{center}
\begin{tcolorbox}[enhanced, breakable, width=0.92\linewidth,
                  colback=gray!3, colframe=gray!45, boxrule=0.5pt,
                  arc=3pt, left=12pt, right=12pt, top=8pt, bottom=8pt,
                  title={Argument flow \quad
                         \disc{v}\,V \;\disc{p}\,P \;\disc{h}\,H \;\disc{c}\,C}]

\sknode{\disc{v}}{Assumption~\ref{ass:cost}}{entry costs a little, and ties split evenly}{%
  $0<\kappa<1/n$ per front opened --- worth opening one, never worth opening all\sktag{setup}}

\skstep{so}

\sknode{\disc{p}}{Proposition~\ref{prop:exist}}{the contest has a mixed-strategy equilibrium}{%
  for any opened fronts, players randomize how hard to fight on each\sktag{engine}}

\skstep{and}

\sknode{\disc{p}}{Lemma~\ref{lem:concentrate}}{each side opens only a few fronts}{%
  at most $\lceil 1/(2\kappa)\rceil$ of them --- a bounded number, not all $n$\sktag{engine}}

\skstep{so}

\sknode{\disc{h}}{Corollary~\ref{cor:sparse}}{those few fronts don't grow with $n$}{%
  the support stays size $O(1/\kappa)$, independent of $n$\sktag{hinge}}

\skstep{so}

\sknode{\disc{c}}{Theorem~\ref{thm:char}}{exact closed-form for how they play}{%
  a symmetric mixed equilibrium: open $m^\star=\lceil 1/(2\kappa)\rceil$ fronts,
  spend uniformly on $[0,2/m^\star]$ each
  \;\textnormal{(needs Proposition~\ref{prop:exist})}\sktag{main result}}

\skstep{still open}

\sknode{\disc{c}}{Uniqueness?}{is this the only solution?}{%
  rule out asymmetric equilibria on the same support --- would turn
  Theorem~\ref{thm:char} green\sktag{open frontier}}

\medskip
\hrule
\smallskip
{\footnotesize\itshape
Read top-to-bottom: each step follows from the ones above; the word on each
$\downarrow$ is the logic. ``(needs $X$)'' flags a step that leans on something
not directly above it. Disc colour = current verification status.}

\end{tcolorbox}
\end{center}


% --- closing verdict: why this rarity rating (1-2 sentences) ---
\bigskip\par\noindent\rule{\linewidth}{0.4pt}\par\smallskip
\noindent{\bfseries Why this rating.}\quad
\rarity{\protorarity}\;(\protorarity/5)\ \textperiodcentered\ {\scshape\prototier}.\par
\smallskip
\noindent\protorarityrationale

\end{document}
